% Generated by scripts/build-paper.py — edit paper/src/survey.src.md
\section{A proposed evaluation protocol}\label{sec:a-proposed-evaluation-protocol}

No experiment is reported in this survey. The protocol below is proposed so that the open questions of §7 can be answered with comparable numbers.

\textbf{Tasks.} Seven compact tracks: intent and semantic classification; pairwise judging; tool and model routing; narrative variable coding; sequential agent control; open-set decisions with the right answer present or absent; and numerically grounded decisions in which the needed quantity must be computed. Visual inputs form a separate track, and text-only systems receive the same extracted features.

\textbf{Systems and baselines.} A majority or rule baseline; TF-IDF or embedding plus a linear classifier; SetFit, ModernBERT and GLiClass \citep{arxiv2209_11055,arxiv2412_13663,arxiv2508_07662}; frozen-LM option readouts with and without rotation averaging; constrained JSON generation with the same backbone \citep{arxiv2411_15100}; open decision models, general and specialised; version-pinned hosted Jev; a strong LLM; and two cascades (cheap trained or generative first stage, Jev first stage). Supervised systems report their training data and cost.

\textbf{Metrics.} Accuracy or macro-F1 with invalid outputs counted as errors; NLL, Brier and ECE with stated binning and bootstrap intervals, keeping primitives separate; coherence residuals for complementary questions and for the same question asked through different interfaces; missing-answer detection and false rejection; risk–coverage curves and coverage at fixed risk, with thresholds and temperatures fitted on validation data only; per-item flip rates under option rotation, name–rubric rebinding, label–definition conflict, neutral and random names and language variants, each against a test–retest floor; attack success for instruction-style and false-state injections; single-request p50/p95, amortized per-question time, throughput, deadline attainment under load and end-to-end time as separate columns; API bills per correct decision, including every pre-screen and escalation call.

\textbf{Splits and records.} Group splits and bootstraps by scenario or document. For each item keep the input, schema, option order, full distribution, model ID and version, timestamp, errors, retries and cost. Register the analysis plan, fix thresholds before held-out evaluation, and repeat a slice on a later day.

\textbf{First experiments.} (i) Name–rubric binding × native calibration × selective escalation on a fixed test set: does low-confidence escalation catch binding-induced errors, and at what cost? (ii) Missing-answer menus × yes/no verification × held-out rejection thresholds. (iii) Coherence of complementary questions × interface × recalibration. Each asks whether structured, probabilistic errors are in fact easier to intercept (F1 × F2 × F3 × F8).
